\thispagestyle{empty}

\begin{comment}

\begin{center}
\vspace*{1.0cm}
    \Large\textbf{Abstract}\\[1.cm]
    
    \normalsize Shor’s algorithm presents a polynomial-time quantum solution to the integer factorization problem, threatening classical encryption protocols like RSA. This article introduces the mathematical and computational foundations required to understand the algorithm, including modular arithmetic, continued fractions, and quantum principles. We provide a step-by-step breakdown of Shor’s algorithm and its quantum subroutine, emphasizing accessibility and conceptual clarity, while also addressing practical constraints in current quantum hardware. This article aims to make the concepts of Shor's algorithm accessible to readers with limited background in mathematics and quantum computing, while maintaining mathematical rigor and conceptual depth throughout.
\end{center}
\end{comment}

\vspace*{1.0cm}

\begin{abstract}
    Shor’s algorithm presents a polynomial-time quantum solution to the integer factorization problem, threatening classical encryption protocols like RSA. This article introduces the mathematical and computational foundations required to understand the algorithm, including modular arithmetic, continued fractions, and quantum principles. We provide a step-by-step breakdown of Shor’s algorithm and its quantum subroutine, emphasizing accessibility and conceptual clarity, while also addressing practical constraints in current quantum hardware. This article aims to make the concepts of Shor's algorithm accessible to readers with limited background in mathematics and quantum computing, while maintaining mathematical rigor and conceptual depth throughout.
\end{abstract}
